% GENERATED FILE. Edit canonical lessons, then run npm run generate:books.
% canonical-source-hash: fnv1a64:fd21c3a977d55554
% canonical-lessons: SA-C551-jahati, SA-C551-anusocati, SA-C551-pascattapah, SA-C551-viraktah, SA-C551-utsrjati

\chapter{Giving Something Up}
\label{ch:sa-giving-something-up}
\hlchaptermodality{551}

\begin{chapteropening}
\textbf{By the end of this chapter:} \emph{I can say what I would give up and what I would regret.}

Complete the last lesson of chapter 551: I can say what I would give up and what I would regret.
\end{chapteropening}

\section[\texorpdfstring{\sk{जहाति} (jahāti)}{जहाति (jahāti)}]{\sk{जहाति} (jahāti) — leaves behind; abandons}

\label{lesson:SA-C551-jahati}



\begin{quote}
Before the new one: say the Sanskrit for puts to use, then the Sanskrit for a handing out.
\end{quote}

\subsection*{You'll want to know: \sk{जहाति}}
\textbf{\sk{जहाति}} — \emph{jahāti} — \textquotedblleft{}leaves behind; abandons\textquotedblright{}.

\begin{quote}
\sk{यदि} \sk{सः} \sk{आलस्यं} \sk{जहाति}, \sk{तर्हि} \sk{सः} \sk{सुखी} \sk{भविष्यति।}
\end{quote}

\begin{quote}
\emph{yadi saḥ ālasyaṃ jahāti, tarhi saḥ sukhī bhaviṣyati.}
\end{quote}

\begin{quote}
If he leaves laziness behind, he will be happy.
\end{quote}

Past (he …): \sk{सः} \sk{गृहं} \sk{हीनवान्}, \emph{saḥ gṛhaṃ hīnavān}.

\begin{cousinweb}
\textbf{\sk{जहाति}} (\emph{jahāti}) comes from the root √hā, \textquotedblleft{}to leave, to abandon\textquotedblright{}, whose present is \emph{jahāti}; its past participle is \emph{hīna} (MW).
\end{cousinweb}

\begin{grammarlens}[title={What you've built}]
One more word for saying what happens if.
\end{grammarlens}

\subsection*{Your turn}
\begin{itemize}
\raggedright
  \item \emph{Say it:} \emph{jahāti}
  \item \emph{Say it:} \emph{jahāti}, once more
  \item \emph{Say it:} \emph{yadi saḥ ālasyaṃ jahāti, tarhi saḥ sukhī bhaviṣyati.} — If he leaves laziness behind, he will be happy.
  \item \emph{Recall:} say the Sanskrit for puts to use, then the Sanskrit for a handing out, then say \emph{jahāti} again
\end{itemize}

\begin{tcolorbox}[breakable,colback=teal!4,colframe=teal!35!black,title={Before you move on}]
What does \sk{जहाति} mean? (\textquotedblleft{}Leaves behind; abandons\textquotedblright{}.) Say it once more.
\end{tcolorbox}
\section[\texorpdfstring{\sk{अनुशोचति} (anuśocati)}{अनुशोचति (anuśocati)}]{\sk{अनुशोचति} (anuśocati) — regrets}

\label{lesson:SA-C551-anusocati}



\begin{quote}
Before the new one: say the Sanskrit for a handing out, then the Sanskrit for leaves behind.
\end{quote}

\subsection*{You'll want to know: \sk{अनुशोचति}}
\textbf{\sk{अनुशोचति}} — \emph{anuśocati} — \textquotedblleft{}regrets\textquotedblright{}.

\begin{quote}
\sk{यदि} \sk{त्वं} \sk{न} \sk{पठसि}, \sk{तर्हि} \sk{पश्चात्} \sk{अनुशोचसि।}
\end{quote}

\begin{quote}
\emph{yadi tvaṃ na paṭhasi, tarhi paścāt anuśocasi.}
\end{quote}

\begin{quote}
If you don't study, you will regret it later.
\end{quote}

Past (he …): \sk{सः} \sk{पश्चात्} \sk{अनुशोचितवान्}, \emph{saḥ paścāt anuśocitavān}.

\begin{cousinweb}
\textbf{\sk{अनुशोचति}} (\emph{anuśocati}) is built on \emph{anu}, \textquotedblleft{}after\textquotedblright{}, and the root √śuc, \textquotedblleft{}to grieve\textquotedblright{}, which also gives \emph{śoka}, \textquotedblleft{}grief\textquotedblright{}; MW gives \emph{anuśuc} as to mourn over, to regret (MW).
\end{cousinweb}

\begin{grammarlens}[title={What you've built}]
One more word for saying what happens if.
\end{grammarlens}

\subsection*{Your turn}
\begin{itemize}
\raggedright
  \item \emph{Say it:} \emph{anuśocati}
  \item \emph{Say it:} \emph{anuśocati}, once more
  \item \emph{Say it:} \emph{yadi tvaṃ na paṭhasi, tarhi paścāt anuśocasi.} — If you don't study, you will regret it later.
  \item \emph{Recall:} say the Sanskrit for a handing out, then the Sanskrit for leaves behind, then say \emph{anuśocati} again
\end{itemize}

\begin{tcolorbox}[breakable,colback=teal!4,colframe=teal!35!black,title={Before you move on}]
What does \sk{अनुशोचति} mean? (\textquotedblleft{}Regrets\textquotedblright{}.) Say it once more.
\end{tcolorbox}
\section[\texorpdfstring{\sk{पश्चात्तापः} (paścāttāpaḥ)}{पश्चात्तापः (paścāttāpaḥ)}]{\sk{पश्चात्तापः} (paścāttāpaḥ) — remorse}

\label{lesson:SA-C551-pascattapah}



\begin{quote}
Before the new one: say the Sanskrit for leaves behind, then the Sanskrit for regrets.
\end{quote}

\subsection*{You'll want to know: \sk{पश्चात्तापः}}
\textbf{\sk{पश्चात्तापः}} — \emph{paścāttāpaḥ} — \textquotedblleft{}remorse\textquotedblright{}.

\begin{quote}
\sk{यदि} \sk{त्वं} \sk{सत्यं} \sk{वदसि}, \sk{तर्हि} \sk{पश्चात्तापः} \sk{न} \sk{भविष्यति।}
\end{quote}

\begin{quote}
\emph{yadi tvaṃ satyaṃ vadasi, tarhi paścāttāpaḥ na bhaviṣyati.}
\end{quote}

\begin{quote}
If you tell the truth, there will be no remorse.
\end{quote}

A masculine noun.

\begin{cousinweb}
\textbf{\sk{पश्चात्तापः}} (\emph{paścāttāpaḥ}) is built from \emph{paścāt}, \textquotedblleft{}after\textquotedblright{}, and \emph{tāpa}, \textquotedblleft{}heat, pain\textquotedblright{}, from the root √tap; MW gives \emph{paścāttāpa} as after-pain, sorrow, regret, repentance (MW).
\end{cousinweb}

\begin{grammarlens}[title={What you've built}]
One more word for saying what happens if.
\end{grammarlens}

\subsection*{Your turn}
\begin{itemize}
\raggedright
  \item \emph{Say it:} \emph{paścāttāpaḥ}
  \item \emph{Say it:} \emph{paścāttāpaḥ}, once more
  \item \emph{Say it:} \emph{yadi tvaṃ satyaṃ vadasi, tarhi paścāttāpaḥ na bhaviṣyati.} — If you tell the truth, there will be no remorse.
  \item \emph{Recall:} say the Sanskrit for leaves behind, then the Sanskrit for regrets, then say \emph{paścāttāpaḥ} again
\end{itemize}

\begin{tcolorbox}[breakable,colback=teal!4,colframe=teal!35!black,title={Before you move on}]
What does \sk{पश्चात्तापः} mean? (\textquotedblleft{}Remorse\textquotedblright{}.) Say it once more.
\end{tcolorbox}
\section[\texorpdfstring{\sk{विरक्तः} (viraktaḥ)}{विरक्तः (viraktaḥ)}]{\sk{विरक्तः} (viraktaḥ) — detached; indifferent}

\label{lesson:SA-C551-viraktah}



\begin{quote}
Before the new one: say the Sanskrit for regrets, then the Sanskrit for remorse.
\end{quote}

\subsection*{You'll want to know: \sk{विरक्तः}}
\textbf{\sk{विरक्तः}} — \emph{viraktaḥ} — \textquotedblleft{}detached; indifferent\textquotedblright{}.

\begin{quote}
\sk{यदि} \sk{सः} \sk{धनात्} \sk{विरक्तः} \sk{अस्ति}, \sk{तर्हि} \sk{सः} \sk{सुखेन} \sk{जीवति।}
\end{quote}

\begin{quote}
\emph{yadi saḥ dhanāt viraktaḥ asti, tarhi saḥ sukhena jīvati.}
\end{quote}

\begin{quote}
If he does not care for wealth, he lives happily.
\end{quote}

\begin{cousinweb}
\textbf{\sk{विरक्तः}} (\emph{viraktaḥ}) is built on \emph{vi}, \textquotedblleft{}away\textquotedblright{}, and the root √rañj, \textquotedblleft{}to be coloured, to be attached\textquotedblright{}, as its past participle; MW gives \emph{virakta} as changed in feeling, averse, indifferent (MW).
\end{cousinweb}

\begin{grammarlens}[title={What you've built}]
One more word for saying what happens if.
\end{grammarlens}

\subsection*{Your turn}
\begin{itemize}
\raggedright
  \item \emph{Say it:} \emph{viraktaḥ}
  \item \emph{Say it:} \emph{viraktaḥ}, once more
  \item \emph{Say it:} \emph{yadi saḥ dhanāt viraktaḥ asti, tarhi saḥ sukhena jīvati.} — If he does not care for wealth, he lives happily.
  \item \emph{Recall:} say the Sanskrit for regrets, then the Sanskrit for remorse, then say \emph{viraktaḥ} again
\end{itemize}

\begin{tcolorbox}[breakable,colback=teal!4,colframe=teal!35!black,title={Before you move on}]
What does \sk{विरक्तः} mean? (\textquotedblleft{}Detached; indifferent\textquotedblright{}.) Say it once more.
\end{tcolorbox}
\section[\texorpdfstring{\sk{उत्सृजति} (utsṛjati)}{उत्सृजति (utsṛjati)}]{\sk{उत्सृजति} (utsṛjati) — lets go of}

\label{lesson:SA-C551-utsrjati}



\begin{quote}
Before the new one: say the Sanskrit for remorse, then the Sanskrit for detached.
\end{quote}

\subsection*{You'll want to know: \sk{उत्सृजति}}
\textbf{\sk{उत्सृजति}} — \emph{utsṛjati} — \textquotedblleft{}lets go of\textquotedblright{}.

\begin{quote}
\sk{यदि} \sk{त्वं} \sk{क्रोधम्} \sk{उत्सृजसि}, \sk{तर्हि} \sk{मनः} \sk{शान्तं} \sk{भविष्यति।}
\end{quote}

\begin{quote}
\emph{yadi tvaṃ krodham utsṛjasi, tarhi manaḥ śāntaṃ bhaviṣyati.}
\end{quote}

\begin{quote}
If you let go of your anger, your mind will be calm.
\end{quote}

Past (he …): \sk{सः} \sk{कपोतम्} \sk{उत्सृष्टवान्}, \emph{saḥ kapotam utsṛṣṭavān}.

\begin{cousinweb}
\textbf{\sk{उत्सृजति}} (\emph{utsṛjati}) is built on \emph{ud}, \textquotedblleft{}out\textquotedblright{}, and the root √sṛj, \textquotedblleft{}to let go, to emit\textquotedblright{}; MW gives \emph{utsṛj} as to let loose, to let go, to set free; its past participle is \emph{utsṛṣṭa} (MW).
\end{cousinweb}

\begin{grammarlens}[title={What you've built}]
One more word for saying what happens if.
\end{grammarlens}

\subsection*{Your turn}
\begin{itemize}
\raggedright
  \item \emph{Say it:} \emph{utsṛjati}
  \item \emph{Say it:} \emph{utsṛjati}, once more
  \item \emph{Say it:} \emph{yadi tvaṃ krodham utsṛjasi, tarhi manaḥ śāntaṃ bhaviṣyati.} — If you let go of your anger, your mind will be calm.
  \item \emph{Recall:} say the Sanskrit for remorse, then the Sanskrit for detached, then say \emph{utsṛjati} again
\end{itemize}

\begin{tcolorbox}[breakable,colback=teal!4,colframe=teal!35!black,title={Before you move on}]
What does \sk{उत्सृजति} mean? (\textquotedblleft{}Lets go of\textquotedblright{}.) Say it once more.
\end{tcolorbox}
